\documentclass[10pt]{article}


\linespread{0.9}
\usepackage[a4paper,margin=0.8in]{geometry}
\usepackage[T1]{fontenc}
\usepackage{lmodern}
\usepackage{amsmath,amssymb,mathtools}
\usepackage{array}
\usepackage{booktabs}
\usepackage{hyperref}
\usepackage{listings}
\usepackage{tikz-cd}

\hypersetup{
  colorlinks=true,
  linkcolor=blue,
  urlcolor=blue,
  citecolor=blue
}

\lstset{
  basicstyle=\ttfamily\small,
  columns=fullflexible,
  keepspaces=true,
  frame=single,
  showstringspaces=false
}

\newcommand{\End}{\mathrm{End}}

\newcommand{\TL}{\mathrm{TL}}
\newcommand{\trm}{\mathrm{tr}^{\mathrm{markov}}}
\newcommand{\Tr}{\mathrm{tr}}
\newcommand{\cl}{\mathrm{cl}}
\newcommand{\qq}[1]{\left[#1\right]_v}

\title{}
\author{}
\date{}

\begin{document}
\section*{\texttt{tl\_jones.py}}
Code for computing Jones polynomials using Temperley--Lieb representations.

\subsection*{Implementation}
The code is organised around the maps
\[
\begin{tikzcd}
B_n
  \arrow[r, "\rho"]
  \arrow[rr, bend right=30, "\rho_\lambda"']
&
\TL_n
  \arrow[r, "\pi_\lambda"]
&
\mathrm{Mat}_{d_\lambda}(\mathbb{Q}(v))
\end{tikzcd}
\]
which we explain now. 

The unital associative $\mathbb{Q}(v)$-algebra \(\TL_n=\TL_n(\delta)\) (where $\delta=v+v^{-1}$) 
has a basis of non-crossing \((n,n)\)-diagrams and the multiplication is stacking diagrams. 
If $a$ and $b$ are two diagrams then $ab$ means place $b$ on top of $a$. 
The algebra has standard generators $1,e_1,\ldots,e_{n-1}$ with $e_i$ having a 
cup between $i$ and $i+1$ and a cap between $i$ and $i+1$. 
The generators satisfy $e_i^2 = \delta e_i$.

The algebra $\TL_n$ is semisimple with decomposition
\[
\TL_n 
\cong \bigoplus_{\substack{\lambda \vdash n \\ \ell(\lambda)\le 2}} 
\End(V_\lambda)
\cong \bigoplus_{\substack{\lambda \vdash n \\ \ell(\lambda)\le 2}} 
\mathrm{Mat}_{d_\lambda}(\mathbb{Q}(v))
\]
which we explain now. The vector space $V_\lambda$ is defined as follows. 
If \(\lambda=(n-r,r)\), then $V_\lambda$ has basis consisting of
non-crossing `link states' on \(n\) points with \(r\) caps and
\(n-2r\) defects. For example, when \(n=5\) and \(\lambda=(3,2)\),
a basis vector of $V_\lambda$ is the diagram
\[
\begin{tikzpicture}[scale=0.9, baseline=-0.2ex]
  \foreach \x in {0,1,2,3,4} {
    \fill (\x,0) circle (1.5pt);
  }
  \draw (0,0) .. controls (0,0.7) and (1,0.7) .. (1,0);
  \draw (3,0) .. controls (3,0.7) and (4,0.7) .. (4,0);
  \draw (2,0) -- (2,1);
\end{tikzpicture}
\]
The algebra $\TL_n$ acts on $V_\lambda$ by stacking. (The paper by Westbury on Temperley--Lieb representations gives a few equivalent ways to describe this representation of $\TL_n$ and a formula for $d_\lambda = \dim(V_\lambda)$.)

For each partition $\lambda = (n-r,r)$ of $n$, the braid group $B_n$ has a `Jones representation' $\rho_\lambda$ given as follows. Define the group homomorphism 
\[
\rho: B_n\to\TL_n^\times,\qquad \rho(\sigma_i)=v^{-1}I-v^{-2}e_i,
\] 
Let \(\pi_\lambda:\TL_n^\times\hookrightarrow\TL_n\to\mathrm{GL}_{d_\lambda}(\mathbb{Q}(v))\) be the projection map. Then the Jones representation $\rho_\lambda$ is 
\[
\rho_\lambda=\pi_\lambda\circ\rho:B_n\to\mathrm{GL}_{d_\lambda}(\mathbb{Q}(v)).
\]

The Jones representations of $B_n$ are useful for the following reason. First note the (unnormalised) Jones polynomial of a braid's closure equals the Markov trace of the braid's image in $\TL_n$:
\[
J_{\cl(\beta)}(v)=\trm_n\bigl(\rho(\beta)\bigr).
\]
This is defined by closing the diagram of $\rho(\beta)$ and sending each closed circle to $\delta=v+v^{-1}$. 

Second, note \(\trm_n\) is a trace on $\TL_n$, so it must equal a linear combination 
of traces on the matrix algebras $\mathrm{Mat}_{d_\lambda}(\mathbb{Q}(v))$. 
But the only traces on $\mathrm{Mat}_{d_\lambda}(\mathbb{Q}(v))$ are 
scalar multiples of the usual matrix trace. Therefore there are scalars \(c_\lambda\) such that
\[
J_{\cl(\beta)}(v)
=\trm_n(\rho(\beta))
=\sum_{\substack{\lambda \vdash n \\ \ell(\lambda)\le 2}}  
c_\lambda\,\Tr(\rho_\lambda(\beta)).
\] 
Writing \(\lambda=(n-r,r)\), we can prove the formula $c_\lambda=\qq{n-2r+1}$.






\newpage





\section{Overview}


\subsection{\texttt{tl\_jones.py}}
\begin{enumerate}
  \item \texttt{quantum\_integer}
  \item \texttt{projlen}
  \item \texttt{gnf\_to\_braid\_word}
  \item \texttt{link\_states}
  \item \texttt{jones\_rep\_braid\_generator}
  \item \texttt{jones\_rep\_braid\_word}
  \item \texttt{jones\_polynomial}
  \item \texttt{\_temperley\_lieb\_matrix}
  \item \texttt{\_apply\_temperley\_lieb\_generator}
  \item \texttt{\_garside\_delta\_word}
  \item \texttt{\_permutation\_to\_braid\_word}
  \item \texttt{\_validate\_gnf}
  \item \texttt{\_inverse\_permutation}
  \item \texttt{\_right\_descent\_set}
  \item \texttt{\_left\_descent\_set}
  \item \texttt{\_validate\_word}
\end{enumerate}




\subsection{\texttt{test\_tl\_jones.py}}
\begin{enumerate}
  \item \texttt{test\_01\_validation\_errors}
  \item \texttt{test\_01a\_gnf\_to\_braid\_word\_examples}
  \item \texttt{test\_01b\_gnf\_to\_braid\_word\_validation\_errors}
  \item \texttt{test\_02\_projlen\_hard\_coded\_cases}
  \item \texttt{test\_03\_link\_states\_examples}
  \item \texttt{test\_04\_link\_state\_dimensions}
  \item \texttt{test\_05\_temperley\_lieb\_relations}
  \item \texttt{test\_06\_braid\_inverse\_relation}
  \item \texttt{test\_07\_braid\_commuting\_relation}
  \item \texttt{test\_08\_braid\_relation}
  \item \texttt{test\_09\_jones\_polynomial\_simple\_cases}
  \item \texttt{test\_10\_jones\_polynomial\_matches\_knotinfo\_random\_sample}
  \item \texttt{load\_knotinfo\_fixture}
\end{enumerate}








\newpage
\section{Function Explanations}

\subsection{\texttt{tl\_jones.py}}

\paragraph{\texttt{quantum\_two}}
Stores the Laurent polynomial \(v+v^{-1}\), which is the quantum integer \([2]_v\).

\paragraph{\texttt{quantum\_integer}}
Builds the quantum integer $[m]_v = v^{m-1} + v^{m-3} + \cdots + v^{-(m-1)}$. 

\paragraph{\texttt{projlen}}
Computes the spread of a Laurent polynomial in the variable \(v\), meaning the difference between the largest and smallest exponent of \(v\) that appears with nonzero coefficient.

\paragraph{\texttt{gnf\_to\_braid\_word}}
Converts classical Garside normal-form data into its braid-word.
It checks that the input is a valid Garside normal form, expands \(\Delta^d\),
and then expands each simple factor into Artin generators.

\paragraph{\texttt{link\_states}}
Constructs the basis of the standard Temperley--Lieb module \(V_{(n-r,r)}\).
A link state is encoded as a tuple: \texttt{-1} means a defect, while a nonnegative entry records the paired endpoint.

\paragraph{\texttt{jones\_rep\_braid\_generator}}
Returns the matrix of the positive braid generator \(\sigma_i\) acting on \(V_{(n-r,r)}\).
It uses the formula $\rho(\sigma_i) = v^{-1}I - v^{-2}E_i$.

\paragraph{\texttt{jones\_rep\_braid\_word}}
Converts a whole braid word into a matrix by multiplying the matrices of its generators.
Positive integers represent positive braid generators and negative integers represent inverses.

\paragraph{\texttt{jones\_polynomial}}
Computes the unnormalised Jones polynomial of the closure of a braid word.
It sums traces over all standard modules: $J_{\mathrm{cl}(\beta)}(v)
= \sum_r [n-2r+1]_v \operatorname{Tr}(\rho_{(n-r,r)}(\beta))$.

\paragraph{\texttt{\_temperley\_lieb\_matrix}}
Builds the matrix of the Temperley--Lieb generator \(E_i\) on the link-state basis.
This is an internal helper used by \texttt{jones\_rep\_braid\_generator}.

\paragraph{\texttt{\_apply\_temperley\_lieb\_generator}}
Applies the local action of \(E_i\) to one single link state.
It handles the four basic cases: adjacent cap, two defects, one defect plus one paired endpoint, and two paired endpoints.

\paragraph{\texttt{\_garside\_delta\_word}}
Builds the chosen Artin-generator word for the Garside half-twist \(\Delta_n\), using the ascending-block convention $[1,2,\dots,n-1,\,1,2,\dots,n-2,\,\dots,\,1]$.

\paragraph{\texttt{\_permutation\_to\_braid\_word}}
Converts a one-line permutation into a positive braid word by a deterministic bubble-left procedure.
This is used to expand each simple factor in the Garside normal form.

\paragraph{\texttt{\_validate\_gnf}}
Checks the full Garside normal-form input for \texttt{gnf\_to\_braid\_word}.
It verifies that each simple factor is a one-line permutation of \(1,\dots,n\),
rejects the identity and full-\(\Delta\) permutations, and checks the classical
Garside descent-set condition $L(\text{next}) \subseteq R(\text{previous})$.

\paragraph{\texttt{\_inverse\_permutation}}
Builds the inverse of a one-line permutation.

\paragraph{\texttt{\_right\_descent\_set}}
Returns the right descent set of a permutation.

\paragraph{\texttt{\_left\_descent\_set}}
Returns the left descent set via the right descent set of the inverse permutation.

\paragraph{\texttt{\_validate\_word}}
Checks that a braid word is an iterable of nonzero integers, checks that each
generator index lies in the range \(1,\dots,n-1\), and converts the word to a tuple.










\newpage
\subsection{\texttt{test\_tl\_jones.py}}

\paragraph{\texttt{test\_01\_validation\_errors}}
Checks that invalid inputs raise \texttt{ValueError} in the expected places.
This is the most basic defensive test in the suite.

\paragraph{\texttt{test\_01a\_gnf\_to\_braid\_word\_examples}}
Checks explicit examples of the Garside-normal-form conversion code, including positive and negative powers of \(\Delta\), individual simple factors, and a short multi-factor example.

\paragraph{\texttt{test\_01b\_gnf\_to\_braid\_word\_validation\_errors}}
Checks that invalid Garside-normal-form inputs are rejected.
This includes bad strand counts, non-integer powers, malformed permutations, forbidden identity or \(\Delta\) factors, and violations of the descent-set rule.

\paragraph{\texttt{test\_02\_projlen\_hard\_coded\_cases}}
Checks \texttt{projlen} on a collection of explicit Laurent polynomials, including monomials, constants, zero, and examples with both positive and negative exponents.

\paragraph{\texttt{test\_03\_link\_states\_examples}}
Checks half a dozen explicit hard-coded link-state bases.
This confirms both the combinatorial construction and the ordering convention of the basis.

\paragraph{\texttt{test\_04\_link\_state\_dimensions}}
Checks that the number of basis states agrees with $\binom{n}{r} - \binom{n}{r-1}$.

\paragraph{\texttt{test\_05\_temperley\_lieb\_relations}}
Checks Temperley--Lieb algebra relations for the internal matrices \(E_i\):
\[
E_i^2 = \delta E_i,\qquad
E_iE_j = E_jE_i \text{ for } |i-j|>1,\qquad
E_iE_{i+1}E_i = E_i.
\]

\paragraph{\texttt{test\_06\_braid\_inverse\_relation}}
Checks the matrix of \(\sigma_i^{-1}\) is the inverse of the one of \(\sigma_i\).

\paragraph{\texttt{test\_07\_braid\_commuting\_relation}}
Checks the braid-group relation $\sigma_i \sigma_j = \sigma_j \sigma_i\ (|i-j|>1)$.

\paragraph{\texttt{test\_08\_braid\_relation}}
Checks the main braid relation $\sigma_i \sigma_{i+1} \sigma_i=\sigma_{i+1} \sigma_i \sigma_{i+1}$.

\paragraph{\texttt{test\_09\_jones\_polynomial\_simple\_cases}}
Checks a small collection of explicit Jones polynomial examples, including the identity braid and several low-strand examples.

\paragraph{\texttt{test\_10\_jones\_polynomial\_matches\_knotinfo\_random\_sample}}
Chooses a (seeded) random sample of 30 knots from the KnotInfo-derived CSV file, prints their names, and checks that the code reproduces the database values after the normalisation change of variables.

\paragraph{\texttt{load\_knotinfo\_fixture}}
Loads the CSV file \texttt{jones\_braid\_knots.csv} and converts each row into a dictionary with parsed braid word and metadata.


\end{document}








